\section{Inertial Parameters Estimation}
To the best of our knowledge, there has been no comprehensive review dedicated specifically to center-of-mass (CoM) estimation for manipulated objects. Existing surveys primarily focus on the identification of the complete set of inertial parameters rather than CoM estimation itself. A survey of Schedlinsk \cite{Schedlinski2001Survey} presents a literature review of existing methodologies on how to estimate these parameters in mechanical engineering, whereas Mavrakis and Stolkin \cite{liu2022inertialsurvey} is specific to robotics. According to Mavrakis and Stolkin's taxonomy, the methods for estimating inertial parameters can be categorized into three main groups by degree of interaction — purely visual, exploratory, and fixed-object.

The first visual methods relied on the assumption of a known and uniform object density in order to simplify the estimation problem by reducing it to the computation of the object's volume. In , \cite{Chien1986Identification} the object is epresented as a collection of equally sized voxels with identical masses, from which the center of mass is determined by computing the moment of inertia matrix. One of the most widely recognized approaches is presented in \cite{Mirtich1996MassProperties}, the author decomposed a rigid body into polyhedral components with uniform density and computed the mass, and inertia tensor of each component before combining them to obtain the inertial properties of the complete object. The computation was achieved by converting the required three-dimensional volume integrals into two-dimensional surface integrals and one-dimensional line integrals using the divergence theorem, projection, and Green's theorem. However, while visual perception can accurately capture an object's shape and volume, the underlying density distribution is generally unknown, may vary throughout the object, and cannot be directly observed solely from visual information.

Exploratory and fixed-object methods, on the other hand, involve physical interaction with the object to perform the estimation task. These interactive techniques are commonly formulated based on the Newton--Euler equations of motion, which establish a mathematical relationship between the forces/torques (wrenches) applied to an object and its resulting motion. Properties like the center of mass (CoM), mass and the inertia tensor are treated as constant parameters within these equations and are subsequently estimated from the resulting system of equations, typically using least-squares techniques. In fixed object methods \cite{Khosla1985Parameter,Swevers2007Dynamic,Gautier1997Dynamic}, the unknown object is rigidly connected to the robotic system, during the estimation process, either through a firm grasp or by being directly attached to the end-effector. Excitation trajectories are then executed by the robotic system to retrieve the necessary data from wrist sensors. The approach proposed in \cite{Dong2018Payload} employed a robotic arm to perform three distinct trajectories to distinguish the individual effects of mass, CoM, and the inertia tensor. The fundamental principle of this method can also be applied to different robotic platforms. In \cite{Mellinger2011Aerial}, the authors use a UAV equipped with two customised grippers to grasp the object and estimated the payload mass, CoM, and inertia from the UAV dynamics and measured motion, while \cite{Lee2017AerialTransportation} extended this concept to a hexacopter equipped with a 2-DOF manipulator, enabling simultaneous payload estimation and adaptive motion control. Although this method can achieve the most accurate f full three-dimensional inertial properties under controlled conditions, its requirements for a secure grasp and specialized force/torque sensors makes it less versatile than exploratory methods in unconstrained or hazardous scenarios.

Exploratory methods exploit the measurements by utilizing more basic interaction such as such as poking, pushing or tilting. As mentioned, early studies of this category similarly estimates the unknown properties by solving the corresponding dynamic equations. In \cite{Yu2005Pushing}, 

